\section*{Capítulo \ref{ch:b1:complex} --- Números complejos}

\begin{solution}{exo:b1:complex:1}
$1 + \iu = \sqrt 2\, \eu^{\iu\pi/4}$;
$\;\sqrt 3 - \iu = 2\, \eu^{-\iu\pi/6}$;
$\;-5 = 5\, \eu^{\iu\pi}$;
\[
\frac{1 + \iu}{\sqrt 3 - \iu}
= \frac{\sqrt 2}{2}\, \eu^{\iu(\pi/4 + \pi/6)}
= \frac{\sqrt 2}{2}\, \eu^{5\iu\pi/12}.
\]
Calculando el mismo cociente por vía algebraica (multiplicando por el \omterm{def:b1:complex:field}{conjugado}):
\[
\frac{(1 + \iu)(\sqrt 3 + \iu)}{4}
= \frac{(\sqrt 3 - 1) + \iu(\sqrt 3 + 1)}{4}.
\]
Identificando con $\frac{\sqrt 2}{2}(\cos\frac{5\pi}{12} +
\iu\sin\frac{5\pi}{12})$:
\[
\cos\frac{5\pi}{12} = \frac{\sqrt 6 - \sqrt 2}{4},
\qquad
\sin\frac{5\pi}{12} = \frac{\sqrt 6 + \sqrt 2}{4}.
\]
\end{solution}

\begin{solution}{exo:b1:complex:2}
$(1+\iu)^2 = 2\iu$, luego $(1+\iu)^{20} = (2\iu)^{10} = 2^{10}\,\iu^{10}
= 1024 \times (-1) = -1024$.

En forma exponencial, $(1+\iu)^n = 2^{n/2}\, \eu^{\iu n\pi/4}$, que es
real si y solo si $\sin\frac{n\pi}{4} = 0$, es decir, $4 \mid n$. Así
pues, $(1+\iu)^n \in \R$ exactamente para los múltiplos de $4$ (con
valor $(-4)^{n/4}$).
\end{solution}

\begin{solution}{exo:b1:complex:3}
$\abs{z - 2} = \abs{z + \iu}$: equidistancia de los puntos $2$ y $-\iu$
--- la mediatriz del segmento que une $(2, 0)$ y $(0, -1)$.

$\abs{z - 1} = 2$: la circunferencia de centro $1$ y radio $2$.

$\frac{z-1}{z+1} \in \iu\R$: por la~\cref{prop:b1:complex:geometry}~(3),
los puntos $M(z)$, $A(1)$, $B(-1)$ cumplen que las rectas $MA$ y $MB$
son perpendiculares (o bien $z = 1$, donde el cociente vale
$0 \in \iu\R$). El \omterm{def:b1:logic:sets}{conjunto} es la circunferencia de diámetro
$[-1, 1]$ (la circunferencia unidad) menos el punto $-1$, donde el
cociente no está definido. \emph{Comprobación por cálculo:}
$z = \eu^{\iu\theta}$ da $\frac{z-1}{z+1} =
\frac{\eu^{\iu\theta/2}(\eu^{\iu\theta/2} -
\eu^{-\iu\theta/2})}{\eu^{\iu\theta/2}(\eu^{\iu\theta/2} +
\eu^{-\iu\theta/2})} = \iu\tan\frac\theta2 \in \iu\R$.
\end{solution}

\begin{solution}{exo:b1:complex:4}
Linealización:
\[
\sin^4\theta
= \Bigl(\frac{\eu^{\iu\theta} - \eu^{-\iu\theta}}{2\iu}\Bigr)^{\!4}
= \frac{\eu^{4\iu\theta} - 4\eu^{2\iu\theta} + 6 - 4\eu^{-2\iu\theta}
+ \eu^{-4\iu\theta}}{16}
= \frac{\cos 4\theta - 4\cos 2\theta + 3}{8}.
\]
Desarrollo: por De Moivre, $\cos 4\theta = \Re\bigl((c +
\iu s)^4\bigr) = c^4 - 6c^2 s^2 + s^4$ con $c = \cos\theta$,
$s = \sin\theta$; sustituyendo $s^2 = 1 - c^2$:
\[
\cos 4\theta = c^4 - 6c^2(1 - c^2) + (1 - c^2)^2
= 8c^4 - 8c^2 + 1 .
\]
\end{solution}

\begin{solution}{exo:b1:complex:5}
$8\iu = 8\,\eu^{\iu\pi/2}$, luego las raíces cúbicas son
$2\,\eu^{\iu(\pi/6 + 2k\pi/3)}$, $k = 0, 1, 2$:
\[
z_0 = 2\eu^{\iu\pi/6} = \sqrt 3 + \iu,\quad
z_1 = 2\eu^{5\iu\pi/6} = -\sqrt 3 + \iu,\quad
z_2 = 2\eu^{3\iu\pi/2} = -2\iu :
\]
un triángulo equilátero sobre la circunferencia de radio $2$.

$-4 = 4\eu^{\iu\pi}$, así que $z^4 = -4$ tiene soluciones $\sqrt 2\,
\eu^{\iu(\pi/4 + k\pi/2)}$: $\;1 + \iu$, $-1 + \iu$, $-1 - \iu$,
$1 - \iu$. Emparejando las raíces \omterm{def:b1:complex:field}{conjugadas}:
\[
X^4 + 4 = \bigl(X^2 - 2X + 2\bigr)\bigl(X^2 + 2X + 2\bigr),
\]
puesto que $(X - (1+\iu))(X - (1-\iu)) = X^2 - 2X + 2$, y análogamente
para el otro par. (Desarróllese para comprobarlo.)
\end{solution}

\begin{solution}{exo:b1:complex:6}
$C_n + \iu S_n = \sum_{k=0}^{n} \eu^{\iu k\theta} =
\frac{\eu^{\iu(n+1)\theta} - 1}{\eu^{\iu\theta} - 1}$ (suma geométrica
de razón $\eu^{\iu\theta} \neq 1$). Factorización por el ángulo mitad
(\cref{met:b1:complex:trig}~(4)) en numerador y denominador:
\[
\frac{2\iu\sin\frac{(n+1)\theta}{2}\, \eu^{\iu(n+1)\theta/2}}
{2\iu\sin\frac{\theta}{2}\, \eu^{\iu\theta/2}}
= \frac{\sin\frac{(n+1)\theta}{2}}{\sin\frac{\theta}{2}}\,
\eu^{\iu n\theta/2}.
\]
Tomando partes real e imaginaria:
\[
C_n = \frac{\sin\frac{(n+1)\theta}{2}}{\sin\frac{\theta}{2}}
\cos\frac{n\theta}{2},
\qquad
S_n = \frac{\sin\frac{(n+1)\theta}{2}}{\sin\frac{\theta}{2}}
\sin\frac{n\theta}{2}.
\]
\end{solution}

\begin{solution}{exo:b1:complex:7}
Discriminante: $\Delta = (3 + 4\iu)^2 - 4(-1 + 5\iu) = 9 + 24\iu - 16 +
4 - 20\iu = -3 + 4\iu$. Raíces cuadradas de $-3 + 4\iu$: resuélvase
$x^2 - y^2 = -3$, $2xy = 4$, $x^2 + y^2 = 5$; entonces $x^2 = 1$,
$y = 2/x$, con los mismos signos: $\delta = \pm(1 + 2\iu)$. Por tanto
\[
z = \frac{(3 + 4\iu) \pm (1 + 2\iu)}{2}
\in \{\, 2 + 3\iu,\; 1 + \iu \,\}.
\]
\emph{Comprobación:} suma $= 3 + 4\iu$ y producto
$(2+3\iu)(1+\iu) = -1 + 5\iu$, como exigen los coeficientes.
\end{solution}

\begin{solution}{exo:b1:complex:8}
\begin{enumerate}
  \item \cref{prop:b1:complex:sumroots} con $n = 5$.
  \item $u + v = \omega + \omega^2 + \omega^3 + \omega^4 = -1$ por (1).
        Para el producto, desarróllese y redúzcanse los exponentes \omterm{def:b1:complex:field}{módulo} $5$:
        \[
        uv = (\omega + \omega^4)(\omega^2 + \omega^3)
        = \omega^3 + \omega^4 + \omega^6 + \omega^7
        = \omega^3 + \omega^4 + \omega + \omega^2 = -1 .
        \]
        Así pues, $u$ y $v$ tienen suma $-1$ y producto $-1$: son las
        dos raíces de $X^2 + X - 1$.
  \item $u = \omega + \conj\omega = 2\cos\frac{2\pi}{5} > 0$ (el ángulo
        es agudo), y la raíz positiva de $X^2 + X - 1$ es
        $\frac{-1 + \sqrt 5}{2}$. Luego $\cos\frac{2\pi}{5} =
        \frac{\sqrt 5 - 1}{4}$.
\end{enumerate}
\end{solution}

\begin{solution}{exo:b1:complex:9}
Desarróllense los dos cuadrados como en la demostración de
la~\cref{prop:b1:complex:rules}~(4):
\[
\abs{z + w}^2 = \abs z^2 + \abs w^2 + 2\Re(z\conj w),
\qquad
\abs{z - w}^2 = \abs z^2 + \abs w^2 - 2\Re(z\conj w),
\]
y súmense. Geométricamente, $\abs{z+w}$ y $\abs{z-w}$ son las
longitudes de las dos diagonales del paralelogramo de vértices
$0, z, z+w, w$, mientras que $\abs z$ y $\abs w$ son las longitudes de
los lados: la suma de los cuadrados de las diagonales es igual a la
suma de los cuadrados de los cuatro lados.
\end{solution}

\begin{solution}{exo:b1:complex:10}
Sea $j = \eu^{2\iu\pi/3}$, de modo que $j^2 + j + 1 = 0$ y
$\eu^{\iu\pi/3} = -j^2$, $\eu^{-\iu\pi/3} = -j$. El triángulo es
equilátero directo si y solo si $c - a = -j^2 (b - a)$, e indirecto si
y solo si $c - a = -j(b - a)$ (\cref{ex:b1:complex:equilateral}).

Considérense $P = a + jb + j^2 c$ y $Q = a + j^2 b + jc$. Usando
$1 + j^2 = -j$ y $j^3 = 1$:
\[
c - a + j^2(b - a) = c + j^2 b - (1 + j^2)\,a = c + j^2 b + ja
= j\,(a + jb + j^2 c) = jP,
\]
de modo que el caso directo se lee $jP = 0 \iff P = 0$; el mismo
cálculo con $j$ en lugar de $j^2$ da $c - a + j(b - a) = j^2 Q$, así
que el caso indirecto se lee $Q = 0$. Por tanto: equilátero (con
cualquier orientación) $\iff PQ = 0$. Desarrollando, con $j + j^2 = -1$:
\[
PQ = a^2 + b^2 + c^2 + (j + j^2)(ab + bc + ca)
= a^2 + b^2 + c^2 - (ab + bc + ca).
\]
Luego el triángulo es equilátero si y solo si
$a^2 + b^2 + c^2 = ab + bc + ca$.
\end{solution}

\begin{solution}{exo:b1:complex:11}
Como los $\omega^k$, $k = 0, \dots, n-1$, son exactamente las $n$
raíces de $X^n - 1$ (\cref{thm:b1:complex:roots}), y el polinomio es
mónico:
\[
X^n - 1 = \prod_{k=0}^{n-1} (X - \omega^k)
= (X - 1) \prod_{k=1}^{n-1} (X - \omega^k).
\]
Dividiendo por $X - 1$: $\;1 + X + \dots + X^{n-1} =
\prod_{k=1}^{n-1} (X - \omega^k)$. Evaluando en $X = 1$ se obtiene $P = n$.

Ahora bien, $1 - \omega^k = -\eu^{\iu k\pi/n}\bigl(\eu^{\iu k\pi/n} -
\eu^{-\iu k\pi/n}\bigr) = -2\iu\,\eu^{\iu k\pi/n} \sin\frac{k\pi}{n}$,
de modo que, tomando \omterm{def:b1:complex:field}{módulos} en $P = n$ (cada $\sin\frac{k\pi}n > 0$
para $1 \leq k \leq n-1$):
\[
n = \abs P = \prod_{k=1}^{n-1} 2\sin\frac{k\pi}{n}
= 2^{n-1} \prod_{k=1}^{n-1} \sin\frac{k\pi}{n},
\qquad\text{luego}\qquad
\prod_{k=1}^{n-1} \sin\frac{k\pi}{n} = \frac{n}{2^{n-1}} .
\]
\end{solution}

\begin{solution}{exo:b1:complex:12}
$z = 1$ no es solución ($2^n \neq 0$), así que la ecuación equivale a
$\bigl(\frac{z+1}{z-1}\bigr)^n = 1$, es decir,
$\frac{z+1}{z-1} = \omega^k$ con $\omega = \eu^{2\iu\pi/n}$ y
$k \in \{0, \dots, n-1\}$. El valor $k = 0$ queda excluido ($z + 1 =
z - 1$ es imposible). Para $1 \leq k \leq n - 1$, resolver
$z + 1 = \omega^k(z - 1)$ da $z(1 - \omega^k) = -1 - \omega^k$, de modo
que, con $\varphi = \frac{2k\pi}{n}$ y las factorizaciones por el
ángulo mitad $1 + \eu^{\iu\varphi} = 2\cos\frac\varphi2\,
\eu^{\iu\varphi/2}$, $\eu^{\iu\varphi} - 1 = 2\iu\sin\frac\varphi2
\,\eu^{\iu\varphi/2}$:
\[
z_k = \frac{1 + \omega^k}{\omega^k - 1}
= \frac{2\cos\frac{k\pi}{n}}{2\iu\,\sin\frac{k\pi}{n}}
= -\iu\,\frac{\cos\frac{k\pi}{n}}{\sin\frac{k\pi}{n}} ,
\]
imaginarios puros, como se afirmaba. La \omterm{def:b1:logic:map}{aplicación}
$t \mapsto \cos t/\sin t$ es \omterm{def:b1:logic:inj}{inyectiva} en $\intoo0\pi$ (es
estrictamente decreciente), luego los $n - 1$ valores $z_k$ son
distintos dos a dos: la ecuación, de grado $n - 1$ una vez desarrollada
(los términos $z^n$ se cancelan), tiene exactamente estas $n - 1$
soluciones.
\end{solution}

\begin{solution}{pb:b1:complex:1}
\textbf{1.} $j = \cos\frac{2\pi}3 + \iu\sin\frac{2\pi}3 = -\frac12
+ \iu\frac{\sqrt3}2$; $j^2 = \eu^{4\iu\pi/3} = -\frac12 -
\iu\frac{\sqrt3}2 = \conj j$; $j^3 = 1$; $1 + j + j^2 = 0$ (suma de
las raíces cúbicas de la unidad, \cref{prop:b1:complex:sumroots});
$j^{-1} = j^2$ (pues $j \cdot j^2 = 1$). Sobre la circunferencia
unidad, $1$, $j$, $j^2$ son los vértices de un triángulo equilátero directo.

\textbf{2.} La rotación de centro $a$ y ángulo $\theta$ fija $a$ y gira
un ángulo $\theta$ todo vector con origen en $a$: $r(z) - a =
\eu^{\iu\theta}(z - a)$, es decir, $r(z) = a + \eu^{\iu\theta}(z - a)$.
Componiendo $r(z) = a + \eu^{\iu\theta}(z-a)$ y $r'(z) = a' +
\eu^{\iu\theta'}(z-a')$:
\[
r' \circ r\,(z) = \eu^{\iu(\theta + \theta')} z +
\text{constante},
\]
una \omterm{def:b1:logic:map}{aplicación} de la forma $Az + B$ con
$A = \eu^{\iu(\theta+\theta')}$ de \omterm{def:b1:complex:field}{módulo} $1$. Por
la~\cref{prop:b1:complex:similitude}, es una rotación de ángulo
$\arg A = \theta + \theta'$ si $A \neq 1$, y una traslación si $A = 1$,
es decir, si $\theta + \theta' \in 2\pi\Z$.

\textbf{3.} $(b, c, p)$ es equilátero si y solo si $\abs{p - b} =
\abs{c - b} = \abs{p - c}$. Escribiendo $q = \frac{p - b}{c - b}$, la
primera igualdad dice $\abs q = 1$ y la segunda $\abs{q - 1} = 1$;
juntas, $q = \eu^{\pm\iu\pi/3}$ (los dos puntos de corte de las
circunferencias $\abs q = 1$ y $\abs{q - 1} = 1$ son
$\frac12 \pm \iu\frac{\sqrt3}2$). Luego
$p = b + \eu^{\pm\iu\pi/3}(c - b)$. Para $a = 0$, $b = 1$, $c = \iu$
(triángulo directo), la elección $\eu^{-\iu\pi/3}$ da
\[
p = 1 + \Bigl(\tfrac12 - \iu\tfrac{\sqrt3}2\Bigr)(\iu - 1)
= \tfrac{1 + \sqrt3}2\,(1 + \iu) \approx 1.37 + 1.37\,\iu ,
\]
que queda al otro lado de la recta $BC$ ($x + y = 1$) respecto de
$a = 0$: hacia fuera. La elección $\eu^{+\iu\pi/3}$ da
$p \approx -0.37 - 0.37\,\iu$, del mismo lado que $a$: hacia dentro.

\textbf{4.} Multiplíquese $P = a + jb + j^2c$ por $j^2$: $j^2 P = j^2 a +
j^3 b + j^4 c = b + jc + j^2 a$ (usando $j^3 = 1$, $j^4 = j$); y por
$j$: $jP = ja + j^2 b + c$. Estas son las dos identidades. Si $P = 0$,
entonces $j^2 P = jP = 0$: el criterio se cumple también para
$(b, c, a)$ y $(c, a, b)$ --- invariancia cíclica (como debe ser: a un
triángulo equilátero le da igual qué vértice se enumere primero).

\textbf{5.} Si $a = b$: $0 = a(1 + j) + j^2 c = -j^2 a + j^2 c$ (usando
$1 + j = -j^2$), luego $c = a$. Si $b = c$: $0 = a + b(j + j^2) = a - b$,
luego $a = b$. Si $a = c$: $0 = a(1 + j^2) + jb = -ja + jb$, luego
$a = b$. En cada caso, los tres puntos coinciden.

\textbf{6.} $(a'z + b') \circ (az + b) = a'a\,z + (a'b + b')$ con
$a'a \neq 0$: la misma forma. La inversa de $z \mapsto az + b$ es
$z \mapsto \frac1a z - \frac ba$, de nuevo de la misma forma. Con la
\omterm{def:b1:logic:map}{aplicación} identidad como elemento neutro, las semejanzas directas
forman un grupo con la composición.

\textbf{7.} $f(z) = az + b$ cumple $f(z_1) = w_1$, $f(z_2) = w_2$ si y
solo si $a z_1 + b = w_1$ y $a z_2 + b = w_2$; restando,
$a(z_2 - z_1) = w_2 - w_1$, de modo que
\[
a = \frac{w_2 - w_1}{z_2 - z_1} \;(\neq 0), \qquad
b = w_1 - a z_1
\]
quedan forzados y, recíprocamente, esta elección sirve: existencia y
unicidad.

\textbf{8.} $\abs{f(z) - f(w)} = \abs{a(z - w)} = \abs a\,\abs{z - w}$:
todas las distancias quedan multiplicadas por $\abs a$. Para la forma:
\[
\frac{f(c') - f(a')}{f(b') - f(a')} = \frac{a(c' -
a')}{a(b' - a')} = \frac{c' - a'}{b' - a'} .
\]
Si dos triángulos $(a', b', c')$ y $(a'', b'', c'')$ tienen la misma
forma, sea $f$ la única semejanza directa con $f(a') = a''$ y
$f(b') = b''$ (pregunta 7); entonces la forma de $(a'', b'', f(c'))$
coincide con la de $(a', b', c')$ y, por tanto, con la de
$(a'', b'', c'')$, y la forma determina el tercer vértice a partir de
los dos primeros: $f(c') = c''$. Recíprocamente, la igualdad de formas
se sigue de la invariancia anterior.

\textbf{9.} $b = f(0) = 1$; $a + b = f(1) = \iu$ da $a = \iu - 1$.
Razón $\abs a = \sqrt2$, ángulo $\arg(\iu - 1) = \frac{3\pi}4$. Punto
fijo:
\[
\zeta = \frac{b}{1 - a} = \frac1{2 - \iu} = \frac{2 + \iu}5 .
\]
Así pues, $f$ es la semejanza directa de centro $\frac{2+\iu}5$, razón
$\sqrt2$ y ángulo $\frac{3\pi}4$.

\textbf{10.} Con $\alpha + \beta = 1$ y $f(z) = az + b$:
\[
f(\alpha a' + \beta b') = a\alpha a' + a\beta b' + b
= \alpha(a a' + b) + \beta(a b' + b)
= \alpha f(a') + \beta f(b') ,
\]
siendo la clave $b = (\alpha + \beta)b$. Los puntos medios
($\alpha = \beta = \frac12$), los baricentros (iterando) y los centros
de Napoleón $\alpha b + \beta c$ de más abajo son, por tanto,
\emph{equivariantes}: transformar el triángulo los transforma en
consecuencia. Así, demostrar un \omterm{def:b1:logic:statement}{enunciado} invariante por semejanza para
un triángulo bien colocado lo demuestra para todos --- la licencia
usada en la pregunta 21.

\textbf{11.} El vértice exterior sobre $[b, c]$ es
$p_a = b + \eu^{-\iu\pi/3}(c - b)$ (pregunta 3), luego el centro es
\[
n_a = \frac{b + c + p_a}3
= \frac{2b + c + \frac{1 - \iu\sqrt3}2\,(c - b)}3
= \frac{(3 + \iu\sqrt3)\,b + (3 - \iu\sqrt3)\,c}6
= \alpha b + \beta c ,
\]
con $\alpha = \frac{3 + \iu\sqrt3}6$ y $\beta = \frac{3 -
\iu\sqrt3}6 = \conj\alpha$. El mismo cálculo sobre los lados $[c, a]$ y
$[a, b]$ da $n_b = \alpha c + \beta a$ y $n_c = \alpha a + \beta b$
(desplazamiento cíclico de los papeles).

\textbf{12.} $\alpha + \beta = \frac{3 + \iu\sqrt3 + 3 -
\iu\sqrt3}6 = 1$. Por tanto
\[
\frac{n_a + n_b + n_c}3
= \frac{(\alpha + \beta)(a + b + c)}3 = \frac{a + b + c}3 :
\]
los dos triángulos comparten baricentro.

\textbf{13.} Agrúpese por $\alpha$ y $\beta$ y apliquénse las
identidades de la pregunta 4 a $P = a + jb + j^2c$:
\[
n_a + j n_b + j^2 n_c
= \alpha\,(b + jc + j^2 a) + \beta\,(c + ja + j^2 b)
= \alpha\,j^2 P + \beta\,j P
= (\alpha j^2 + \beta j)\,P .
\]

\textbf{14.} Con $j = \frac{-1 + \iu\sqrt3}2$:
\[
\alpha j = \frac{(3 + \iu\sqrt3)(-1 + \iu\sqrt3)}{12}
= \frac{-3 + 3\iu\sqrt3 - \iu\sqrt3 - 3}{12}
= \frac{-3 + \iu\sqrt3}6 = -\beta ,
\]
de modo que $\alpha j + \beta = 0$, luego $\alpha j^2 + \beta j =
j(\alpha j + \beta) = 0$, y la pregunta 13 da $n_a + jn_b + j^2n_c = 0$
para todo triángulo $(a, b, c)$. Por el criterio (preguntas 4--5), los
centros forman un triángulo equilátero directo o un solo punto: el
teorema de Napoleón.

\textbf{15.} El vértice interior es $b + \eu^{+\iu\pi/3}(c - b)$, y el
cálculo de la pregunta 11 con $\eu^{+\iu\pi/3} = \frac{1 +
\iu\sqrt3}2$ intercambia $\alpha$ y $\beta$: $m_a = \beta b + \alpha
c$, $m_b = \beta c + \alpha a$, $m_c = \beta a + \alpha b$. Usando las
identidades análogas $b + j^2c + ja = j\,Q$ y $c + j^2a + jb = j^2 Q$
para $Q = a + j^2b + jc$:
\[
m_a + j^2 m_b + j m_c
= \beta(b + j^2 c + ja) + \alpha(c + j^2 a + jb)
= (\beta j + \alpha j^2)\, Q = j(\beta + \alpha j)\,Q = 0 ,
\]
puesto que $\alpha j = -\beta$ (pregunta 14). Así pues, los centros
interiores cumplen el criterio de equilátero \emph{indirecto}:
equilátero con la orientación opuesta (o un punto).

\textbf{16.} Por la pregunta 5 aplicada a la terna $(n_a, n_b, n_c)$
(que cumple el criterio directo), dos centros coinciden si y solo si
coinciden los tres; y los tres coinciden si y solo si se cumplen
\emph{los dos} criterios, es decir, si además $n_a + j^2 n_b + j n_c =
0$. Calcúlese como en la pregunta 13, con las identidades
$b + j^2c + ja = jQ$, $c + j^2a + jb = j^2Q$:
\[
n_a + j^2 n_b + j n_c
= \alpha\,jQ + \beta\,j^2 Q = j(\alpha + \beta j)\,Q .
\]
Directamente: $\beta j = \frac{(3 -
\iu\sqrt3)(-1 + \iu\sqrt3)}{12} = \frac{-3 + 3\iu\sqrt3 +
\iu\sqrt3 + 3}{12} = \frac{\iu\sqrt3}3$, luego $\alpha + \beta j =
\frac{3 + \iu\sqrt3 + 2\iu\sqrt3}6 = \frac{1 + \iu\sqrt3}2 \neq 0$. Por
tanto, el triángulo de Napoleón exterior degenera si y solo si
$Q = a + j^2b + jc = 0$, es decir, si y solo si $(a, b, c)$ es un
triángulo equilátero indirecto --- en ese caso las construcciones
«hacia fuera» apuntan todas hacia dentro y comparten un mismo centro.

\textbf{17.} Con $a = 0$, $b = 1$, $c = \iu$:
\[
n_a = \alpha + \beta\iu = \frac{(3 + \sqrt3)(1 + \iu)}6,
\qquad
n_b = \alpha\iu = \frac{-\sqrt3 + 3\iu}6,
\qquad
n_c = \beta = \frac{3 - \iu\sqrt3}6 .
\]
Lados al cuadrado: $n_a - n_b = \frac{(3 + 2\sqrt3) +
\iu\sqrt3}6$ da $\abs{n_a - n_b}^2 = \frac{(3 + 2\sqrt3)^2 +
3}{36} = \frac{24 + 12\sqrt3}{36} = \frac{2 + \sqrt3}3$;
$n_b - n_c = \frac{(3 + \sqrt3)(-1 + \iu)}6$ da $\abs{n_b - n_c}^2 =
\frac{2(3 + \sqrt3)^2}{36} = \frac{24 + 12\sqrt3}{36}$;
$n_c - n_a = \frac{-\sqrt3 - \iu(3 + 2\sqrt3)}6$ da otra vez el mismo
valor. Los tres lados al cuadrado valen $\frac{2 + \sqrt3}3$:
equilátero, como se prometía.

\textbf{18.} Desarróllese:
\[
(a - b)(c - d) + (a - d)(b - c)
= (ac - ad - bc + bd) + (ab - ac - bd + cd)
= ab - ad - bc + cd ,
\]
y $(a - c)(b - d) = ab - ad - bc + cd$: iguales.

\textbf{19.} Tómense \omterm{def:b1:complex:field}{módulos} en la pregunta 18 y aplíquese la
desigualdad triangular (\cref{prop:b1:complex:rules}~(4)):
\[
AC \cdot BD = \abs{(a-b)(c-d) + (a-d)(b-c)}
\leq \abs{a-b}\,\abs{c-d} + \abs{a-d}\,\abs{b-c}
= AB \cdot CD + AD \cdot BC ,
\]
con igualdad si y solo si los dos sumandos están en una misma
semirrecta desde $0$ (el caso de igualdad de la desigualdad triangular).

\textbf{20.} Factorización por el ángulo mitad
(\cref{met:b1:complex:trig}~(4)):
$\eu^{\iu\theta} - \eu^{\iu\varphi} = 2\iu\,\sin\frac{\theta -
\varphi}2\;\eu^{\iu(\theta + \varphi)/2}$, y tomar \omterm{def:b1:complex:field}{módulos} elimina los
factores unimodulares: $\abs{\eu^{\iu\theta} - \eu^{\iu\varphi}}
= 2\,\abs{\sin\frac{\theta - \varphi}2}$.

\textbf{21.} Con $p = \eu^{\iu\theta}$, $\theta \in
\intoo{2\pi/3}{4\pi/3}$, la pregunta 20 da
\[
PA = 2\,\abs{\sin\tfrac\theta2},
\quad
PB = 2\,\abs{\sin\bigl(\tfrac\theta2 - \tfrac\pi3\bigr)},
\quad
PC = 2\,\abs{\sin\bigl(\tfrac\theta2 - \tfrac{2\pi}3\bigr)} .
\]
Sobre el arco, $\frac\theta2 \in \intoo{\pi/3}{2\pi/3}$: entonces
$\sin\frac\theta2 > 0$; $\frac\theta2 - \frac\pi3 \in
\intoo0{\pi/3}$, luego el segundo seno es positivo; $\frac\theta2 -
\frac{2\pi}3 \in \intoo{-\pi/3}0$, luego el tercero es negativo y
$PC = 2\sin\bigl(\frac{2\pi}3 - \frac\theta2\bigr)$. Transformando la
suma en producto:
\[
\sin\Bigl(\frac\theta2 - \frac\pi3\Bigr) +
\sin\Bigl(\frac{2\pi}3 - \frac\theta2\Bigr)
= 2\,\sin\frac\pi6\,\cos\Bigl(\frac\theta2 - \frac\pi2\Bigr)
= \sin\frac\theta2 ,
\]
luego $PB + PC = 2\sin\frac\theta2 = PA$: el teorema de Van Schooten.

\textbf{22.} En $p = -1$: $PA = 2$, $PB = PC = 1$, y todos los lados
del triángulo equilátero miden $\sqrt3$, de modo que los dos miembros
de Ptolomeo valen $2\sqrt3$. Algebraicamente, usando $-1 - j^2 = j$ y
$1 + j = -j^2$:
\[
(a - b)(p - c) = (1 - j)\,(-1 - j^2) = (1 - j)j = j - j^2
= \iu\sqrt3 ,
\]
\[
(a - c)(b - p) = (1 - j^2)(j + 1) = 1 + j - j^2 - j^3
= j - j^2 = \iu\sqrt3 .
\]
Los dos términos son iguales, luego su cociente es $1 \in \R_{>0}$: el
caso de igualdad de la pregunta 19, que reproduce exactamente
$PA \cdot BC = PB \cdot AC + PC \cdot AB$.

\textbf{23.} Para $(a, b, c) = (1, j, j^2)$: $n_a = \alpha j +
\beta j^2 = -\beta + \beta j^2$ (pregunta 14) $= \beta(j^2 - 1)$;
numéricamente, $\beta(j^2 - 1) = \frac{(3 - \iu\sqrt3)}6 \cdot
\bigl(-\frac32 - \iu\frac{\sqrt3}2\bigr) = -1$. Análogamente,
$n_b = \alpha j^2 + \beta = -j$ y $n_c = \alpha + \beta j = -j^2$. El
triángulo de Napoleón exterior de $(1, j, j^2)$ es $(-1, -j, -j^2)$: el
triángulo de partida reflejado a través de su baricentro $0$ --- del
mismo tamaño, girado media vuelta. Un triángulo equilátero es una forma
fija de la construcción de Napoleón, no un límite que se encoge.

\textbf{24.} (i) La multiplicación por un número unimodular como
rotación construyó los vértices y los centros (preguntas 2--3 y 11) y
convirtió distancias sobre la circunferencia en senos (pregunta 20).
(ii) La estructura de grupo y el invariante de forma justificaron
normalizar la circunferencia circunscrita a la circunferencia unidad y
el triángulo a $(1, j, j^2)$ en la pregunta 21, mediante la
equivariancia de la pregunta 10. (iii) La factorización por el ángulo
mitad movió tanto la pregunta 20 como el paso de suma a producto que
remata Van Schooten.

\textbf{25.} El diccionario convierte los \omterm{def:b1:logic:statement}{enunciados} geométricos en
identidades polinómicas en las que toda hipótesis es una ecuación: lo
que aporta es mecanización --- el teorema de Napoleón reducido a
$\alpha j + \beta = 0$, una línea de aritmética en $\Q(\iu\sqrt3)$ ---;
lo que cuesta es visibilidad geométrica --- el cálculo certifica, pero
no muestra \emph{por qué} los centros se cierran en un triángulo
equilátero. Una respuesta razonable es que la identidad explica la
\emph{inevitabilidad} del teorema (se cumple idénticamente en
$a, b, c$, de modo que no interviene ninguna astucia de configuración),
mientras que el dibujo explica su \emph{contenido}. El mismo
diccionario vuelve en forma matricial cuando las rotaciones se
convierten en las matrices ortogonales $2 \times 2$
de los~\cref{ch:b1:matrices} y \cref{ch:b1:euclid}, donde «multiplicar
por $\eu^{\iu\theta}$» es el prototipo de isometría lineal.
\end{solution}
